% =========================
% Chapter 1
% =========================
\chapter{Introduction and Motivation}
\label{ch:introduction}

\section{Motivation}
\label{sec:motivation}

Visual impairment affects approximately 2.2 billion people worldwide, including 39
million who are completely blind~\cite{bourne2017global}. These numbers continue
rising due to population growth and aging, making independent navigation a critical
challenge for millions of individuals.

Among the most pressing difficulties is indoor navigation, which remains largely
unsolved despite significant advances in outdoor GPS-based
systems~\cite{fernandes2019review}. The inability to navigate independently in
unfamiliar indoor spaces --- offices, hospitals, shopping centers --- severely limits
employment opportunities, social participation, and overall quality of life for blind
individuals.

Traditional mobility aids have inherent limitations. The white cane detects obstacles
only within physical reach, missing overhead hazards and providing no advance
warning~\cite{dakopoulos2010wearable}. Guide dogs, while effective, require extensive
training, cost between \$15{,}000 and \$50{,}000, and are not suitable for all users
due to various constraints~\cite{kandalan2020techniques}. Indoor environments further
complicate navigation as building structures attenuate GPS signals, rendering them
unreliable~\cite{li2018vision}.

Recent electronic travel aids have attempted to address these challenges through
continuous scene narration --- constantly verbalizing environmental information.
However, this approach introduces a critical problem: cognitive overload. Users must
process continuous auditory streams while simultaneously attending to ambient sounds,
tactile feedback, and spatial reasoning~\cite{martinez2014cognitive}. Research shows
this can be overwhelming and counterproductive, causing users to miss critical
details.

This tension between providing sufficient navigation information and avoiding
cognitive overload motivates a different approach. Rather than continuous
environmental description, a system can activate perception only when the user
explicitly requests assistance for a specific task, guide that task to completion,
and then fall silent. This task-oriented paradigm reduces cognitive load, preserves
user control, and matches how blind individuals actually operate indoors: they do not
want a running commentary --- they want to \emph{find the cup}, \emph{reach the
kitchen}, or \emph{grasp the object in front of them}.

\section{Problem Statement}
\label{sec:problem_statement}

Current assistive navigation systems for blind users typically fall into two
categories, each with fundamental limitations. The first category comprises
continuous narration systems that constantly describe the surrounding environment
through audio feedback. These systems attempt to provide comprehensive scene
understanding by continuously verbalizing detected objects, spatial relationships,
and environmental changes. While this approach appears thorough, it creates severe
cognitive load problems: users must simultaneously process continuous auditory
information while attending to ambient sounds, tactile feedback from their mobility
aids, and spatial reasoning~\cite{hossain2010cognitive}. Research on cognitive load
in assistive technologies has demonstrated that this constant information flow can
overwhelm users, leading to reduced task performance and increased mental fatigue.

The second category consists of query-based systems where users must explicitly
request information about specific aspects of their
environment~\cite{li2018vision, khan2020insight}. While these systems reduce
cognitive load by providing information only on demand, they introduce a different
problem: users must know precisely what to ask for and when. In unfamiliar
environments, blind users often cannot anticipate what information they need until it
becomes critically necessary. This creates a reactive rather than proactive
assistance paradigm, where users may miss important environmental cues simply because
they did not know to query for them.

Both approaches share a common limitation: they treat perception as either
continuously active or completely passive. Continuous systems waste computational
resources and overwhelm users with unnecessary information; query-based systems place
the burden of managing perception entirely on the user. Neither aligns with how blind
individuals naturally navigate --- performing specific tasks such as locating
objects, reaching destinations, or avoiding obstacles, each with distinct information
requirements and a clear beginning and end.

What remains missing is a middle ground: a system that activates perception and
provides guidance in response to explicit user-initiated tasks, but does so
intelligently without requiring users to specify every detail. The core research
challenge is designing an interaction paradigm and system architecture that enables
\emph{task-scoped perception} --- where visual processing and guidance activate only
during explicitly requested tasks, provide structured assistance toward task
completion, and return to an idle state once the task concludes. A second challenge
follows directly from it: for the hardest task family, indoor navigation, the system
must guide a user to a destination room \emph{without maps, localization
infrastructure, or prior knowledge of the building} --- using only what a phone
camera sees and what the user can naturally contribute. This work addresses the
research question: \textit{How can assistive navigation systems provide reliable,
context-aware guidance while minimizing cognitive load through task-oriented
perception activation?}

\section{High-Level Methodology}
\label{sec:high_level_methodology}

To address the gap identified above, Lumen introduces a task-oriented assistive
vision system that activates perception only in response to explicit user commands.
The user opens a web page on an ordinary smartphone; no dedicated hardware and no
application install are required. The phone's browser captures camera frames,
push-to-talk voice commands, and compass headings, and streams them over a single
WebSocket connection to a backend server that performs all computation: speech
recognition, command parsing, object detection, hand tracking, spatial reasoning, and
speech synthesis. Guidance returns to the user as synthesized voice, complemented by
short haptic cues.

The initial project proposal envisioned a purpose-built wearable device. During
development this was deliberately replaced by the smartphone-browser architecture:
the phone already integrates the camera, microphone, speaker, compass, and wireless
connectivity that the wearable would have had to replicate, and serving the client as
a web page eliminates manufacturing, installation, and platform-approval barriers
entirely (Section~\ref{sec:design_evolution} details this evolution). The
architectural principle is unchanged: a thin sensing client, with all intelligence in
a server the client reaches over the network.

The system organizes assistance into two task families:

\begin{itemize}
    \item \textbf{Object Allocation}: locates a requested item (``find my cup'') and
          guides the user toward it with direction and distance cues derived from
          per-class size priors. Once the object is within arm's reach, the task
          enters its \emph{reach guidance} phase: it tracks the user's hand and
          steers it to the target (``move your hand to the right \ldots{} almost
          there''), completing the task automatically when the fingertip reaches
          the object.
    \item \textbf{Navigation}: guides the user to a destination room
          (``navigate to the kitchen'') through \emph{goal-directed exploration}:
          the system scans each room with a compass-anchored $360^\circ$ sweep,
          recognizes arrival from room-indicator objects (a refrigerator and an oven
          imply a kitchen), and otherwise walks the user to a detected door and
          through it, room by room, warning about obstacles along the walking path.
\end{itemize}

A finite state machine (FSM) manages the system's operational states, ensuring
perception activates only during explicit user tasks. Every transition is driven by
an explicit event --- a voice command, a recognized completion, a cancellation, or a
detected success --- making system behavior predictable, interpretable, and testable.

\section{Contributions}
\label{sec:contributions}

The main contributions of this project are:

\begin{itemize}
    \item \textbf{Task-Scoped Perception Paradigm}: a task-oriented perception
          paradigm that activates vision processing only during explicit user
          commands, eliminating continuous cognitive load while maintaining
          proactive assistance within each task.

    \item \textbf{Goal-Directed Exploration Navigation Without Localization}: an
          indoor navigation approach that requires no maps, markers, or SLAM. The
          user names only the destination; the system explores door-by-door,
          detects arrival semantically from room-indicator objects, and treats the
          user as a direction oracle only when perception is genuinely ambiguous.

    \item \textbf{A Safety-Hardened Door Perception Funnel}: a three-layer door
          detection pipeline --- a custom-trained single-class door detector
          (mAP@50 $=0.946$), geometric gates against walls and whole-scene false
          positives, and a second 4-class model as a semantic verifier --- designed
          failure-by-failure against real phone footage, because a false door
          points a blind user at a wall.

    \item \textbf{Zero-Install Commodity-Hardware Delivery}: the complete system
          runs from a phone browser against a single Python server over one
          WebSocket, with an iOS-compatible audio pipeline and tunnel-based HTTPS
          deployment in minutes.

    \item \textbf{FSM-Governed, Testable Assistive Architecture}: explicit
          state-machine control of every session, with the decision logic of both
          task families isolated from models and I/O and verified by 281
          automated tests, including scripted multi-room navigation journeys.
\end{itemize}

\section{Scope}
\label{sec:scope}

Lumen targets indoor residential and office-like environments: rooms connected by
doors, containing everyday objects from the COCO vocabulary. Object Allocation
supports a curated set of common object classes; Navigation supports six destination
room types recognized by their characteristic objects (kitchen, bathroom, bedroom,
living room, dining room, office). The system complements --- and explicitly does not
replace --- the white cane: its value is advance information in the 2--4\,m range,
not real-time collision avoidance. English voice interaction is supported. The
demonstrated deployment serves one active user per task session, with the
architecture designed for per-session isolation.

\section{Report Outline}
\label{sec:report_outline}

The remainder of this report is organized as follows.
\textbf{Chapter~\ref{ch:background}} provides the foundational concepts behind the
system: object detection, speech recognition and synthesis, language understanding,
and state-based control. \textbf{Chapter~\ref{ch:related_work}} surveys assistive
vision systems, interaction paradigms, indoor navigation approaches, and control
architectures, positioning Lumen within the research landscape.
\textbf{Chapter~\ref{ch:system_design}} presents the system architecture: the
client--server design, the wire protocol, the session model, and the finite state
machine, including the design's evolution from the original wearable proposal.
\textbf{Chapter~\ref{ch:task_families}} details the two task families and their
algorithms --- Object Allocation, Reach Guidance, and goal-directed exploration
Navigation, including the door detector training and the obstacle watchdog.
\textbf{Chapter~\ref{ch:implementation}} describes the implementation: the technology
stack, the per-session task engines, the speech pipeline, deployment, and the
automated test infrastructure. \textbf{Chapter~\ref{ch:evaluation}} evaluates the
system through model metrics, automated behavioral verification, and live trials, and
discusses limitations. \textbf{Chapter~\ref{ch:conclusion}} concludes and outlines
future work.